
% \documentclass[twocolumn]{article}
\documentclass{article}
\usepackage{graphicx,amsmath,fancyhdr,fancyvrb,bm,hyperref}




\hypersetup{
	colorlinks=true,
	linkcolor=blue,
	filecolor=magenta,      
	urlcolor=cyan,
	pdftitle={Overleaf Example},
	pdfpagemode=FullScreen,
}

\urlstyle{same}

\title{Acoustic simulation documentation}
\author{...}
\date{2026}
\pagenumbering{arabic}     % Use Arabic numerals (1, 2, 3, ...)
\setcounter{page}{0}       % Start from page number 0
\thispagestyle{empty}      % Removes number on the current page (the first one)
\usepackage[a4paper,
            bindingoffset=0in,
            left=1.5cm,
            right=1.5cm,
            top=2cm,
            bottom=1.5cm,
            footskip=.25in]{geometry}
\pagestyle{fancy}
\fancyhf{} % clear all header and footer fields
% Put section title on the top (left or right depending on odd/even pages)
\fancyhead[L]{\nouppercase{\leftmark}}
% Page number at the bottom center
\fancyfoot[C]{\thepage}
% invert colors comment untill star to get a printable version
\usepackage{xcolor}
\definecolor{mybg}{RGB}{30,30,30}
\usepackage{background}
\backgroundsetup{
  scale=1,
  % color=black,
  % color=darkgray,
  % color=mybg,
  color=white,
  opacity=1,
  angle=0,
  position=current page.center, % center aligns full-page rectangle
  contents={\rule{\paperwidth}{\paperheight}}
}
% \color{white}
% star

% ----------------------------
% ---------- MACROS ----------
% ----------------------------
% numberless section title + add this to the header
\newcommand{\ms}[1]{%
  \newpage%
  \section*{Week #1, Lecture #1}%
  \markboth{Week #1, Lecture #1}{}%
  \setcounter{equation}{0}
  \setcounter{section}{#1}
  \setcounter{tocdepth}{1}
  \addcontentsline{toc}{section}{Week #1, Lecture #1}
}
% partial derivative, fraction
\newcommand{\pd}[2]{\frac{\partial #1}{\partial #2}}
% upper dagger
\newcommand{\adj}[0]{^\dag}
% braket, bra, ket notations
\newcommand{\braket}[2]{\left\langle #1 \middle| #2 \right\rangle}
\newcommand{\ketbra}[2]{\left| #1 \middle\rangle \middle\langle #2 \right|}
\newcommand{\bra}[1]{\left\langle #1 \right\vert}
\newcommand{\ket}[1]{\left\vert #1 \right\rangle}
% scalar product
\newcommand{\scp}[2]{\left\langle #1 , #2 \right\rangle}
% better sum
\newcommand{\summ}[2]{\sum_{#1}^{#2}}
% evaluated at
\newcommand{\at}[2]{\left. #1 \right\rvert_{#2}}
% array vertical column form stack exchange
\makeatletter
\renewcommand*\env@matrix[1][*\c@MaxMatrixCols c]{%
  \hskip -\arraycolsep
  \let\@ifnextchar\new@ifnextchar
  \array{#1}}
\makeatother
% set,where { #1 | #2 }
\newcommand{\setwhere}[2]{\begin{Bmatrix}[c|c] #1 & #2 \end{Bmatrix}}
% norm
\newcommand{\norm}[1]{\left\vert \left\vert #1 \right\vert \right\vert}
% reference equation
\newcommand{\refeq}[1]{equation (\ref{ #1 })}
% Hilbert-Schmidt scalar product
\newcommand{\scphs}[2]{\scp{#1}{#2}_{H-Sch}}
% quantum system state space (S^+_1(\mathcal{H}))
\newcommand{\qss}[0]{S^+_1(\mathcal{H})}
\newcommand{\qssi}[0]{S^+_1(\mathcal{H}_I)}
\newcommand{\qssii}[0]{S^+_1(\mathcal{H}_{II})}
\newcommand{\qssprod}[0]{S^+_1(\mathcal{H}_I\otimes\mathcal{H}_{II})}
\newcommand{\qssv}[0]{S(\mathcal{H})}
\newcommand{\qssiv}[0]{S(\mathcal{H}_I)}
\newcommand{\qssiiv}[0]{S(\mathcal{H}_{II})}
\newcommand{\qssprodv}[0]{S(\mathcal{H}_I\otimes\mathcal{H}_{II})}
% quantum system state space (S^+_1(  ))
\newcommand{\qssf}[1]{S^+_1(#1)}
% align* equation numbering
\newcommand{\nt}{\stepcounter{equation}\tag{\theequation}}
% in equation, often used words
\newcommand{\und}[0]{\text{ and }}
\newcommand{\oder}[0]{\text{ or }}
\newcommand{\tr}[1]{\text{ Tr}\left( #1 \right)}
\newcommand{\Span}[1]{\text{ Span}\left\{ #1 \right\}}
% \newcommand{\span}[1]{\text{ Span }\left( #1 \right) }
\newcommand{\carray}[1]{\begin{pmatrix} #1 \end{pmatrix}}
\newcommand{\lin}[1]{\text{Lin}\left( #1 \right)}
\newcommand{\subsectiontoc}[1]{%
  \subsection*{#1}%
  \addcontentsline{toc}{subsection}{#1}%
}
\newcommand{\subsubsectiontoc}[1]{%
  \subsubsection*{#1}%
  \addcontentsline{toc}{subsubsection}{#1}%
}
\newcommand{\paragraphtoc}[1]{%
  \paragraph{#1}%
  \addcontentsline{toc}{paragraph}{#1}%
}

\setlength{\parindent}{0pt} % removes paragraph indentation

% equation numbering as <#section.x>
\numberwithin{equation}{section}
% ------------------------------
% ---------- DOCUMENT ----------
% ------------------------------
\begin{document}

\maketitle
\thispagestyle{empty}      % Removes number on the current page (the first one)

\newpage

\setcounter{tocdepth}{4}
\tableofcontents  % <-- generates the table of contents here
\newpage           % optional, starts content on a new page


\section{the pde}

derivation at the wikipedia (TODO: provide link) \\
\url{https://en.wikipedia.org/wiki/Acoustic\_wave\_equation}

\begin{align}
	\pd{^2p(\bm x,t)}{t^2} = a^2 \Delta p(\bm x,t) = a^2 \left( \pd{^2p}{x^2} + \pd{^2p}{y^2} \right)
\end{align}

here $\Delta$ is the Laplace operator

\section{the boundary conditions}

Homogeneous Dirichlet: $\at{p(\bm x,t)}{\bm x \in S_d}=0$

Homogeneous Neumann: $\at{\nabla_{\bm x}p(\bm x,t)\cdot\bm n_n}{\bm x \in S_n} = 0$ % $\pd{p(\bm x_n,t)}{\bm x}=0$

Inhomogeneous Neumann (Transducer): $\at{\nabla_{\bm x}p(\bm x,t)\cdot\bm n_t}{\bm x \in S_t} = \rho\omega^2d\sin(\omega t)$ % $\pd{p(\bm x_t,t)}{\bm x}=\rho\omega^2d\sin(\omega t)$ 

(TODO: is this correct for the transducer?????)
(TODO: add derivation later)

Where $\bm n_*$ is the normal direction of the boundary\\
adn $S_*$ is the set of given boundary points \\
$\rho$ is the fluid density\\
$\omega$ is the excitations angular frequency from the transducer\\
$d$ is the displacement of the transducer

\section{discretization}

notations:
\begin{align}
	\pd{f(\bm x, t)}{t} = \dot f\\
	\pd{g(x,t)}{x} = g'\\
\end{align}
for a uniform grid of:\\
$(t_0,t_1,...,t_n,...,t_N)$\\
$(x_0,x_1,...,x_j,...,x_J)$\\
$(y_0,y_1,...,y_j,...,y_J)$\\
the following notation is adapted: $p(x_j,y_i,t_n) = p^n_{j,i}$, and for 1D: $p(x_j,t_n) = p^n_{j}$\\
here $x_0,y_0,x_J,y_J$, are at the boundary and $x_J-x_0=y_J-y_0=L$, the length of the square simulation domain\\
$\Delta t = t_{n+1} - t_n$\\
$\Delta x = x_{j+1} - x_j = \Delta y = y_{i+1} - y_i$

\subsection{CFL}
\begin{align}
	\sigma = \frac{a \Delta t}{\Delta x} < 1 
\end{align}

\begin{align}
	\Delta t \text{ proportional }  \Delta x \Rightarrow O(\Delta t) = O(\Delta x)
\end{align}

\subsection{pde}

the central difference method is used, that can be seen at $\ref{eq:cdm}$

\begin{align}\label{eq:cdm}
	\left. \pd{^2f}{x^2}\right|_{x=x_j} \approx \frac{f_{j+1} - 2 f_{j} + f_{j-1}}{\Delta x^2} + O(\Delta x^2)
\end{align}

\begin{align}\label{eq:disc}
	\frac{p^{n+1}_{j,i} - 2 p^{n}_{j,i} + p^{n-1}_{j,i}}{\Delta t^2} = a^2\left( \frac{p^{n}_{j+1,i} - 2 p^{n}_{j,i} + p^{n}_{j-1,i}}{\Delta x^2} + \frac{p^{n}_{j,i+1} - 2 p^{n}_{j,i} + p^{n}_{j,i-1}}{\Delta y^2}\right) + 
\end{align}

\begin{align}\label{eq:disc}
	p^{n+1}_{j,i} = 2 p^{n}_{j,i} - p^{n-1}_{j,i} + \frac{a^2\Delta t^2}{\Delta x^2}\left( p^{n}_{j+1,i} + p^{n}_{j-1,i} + p^{n}_{j,i+1} + p^{n}_{j,i-1} - 4 p^{n}_{j,i}\right)
\end{align}

\subsection{bc}
\subsubsection{Homogeneous Dirichlet in 1D}
\begin{align}
	p^n_{j_d} = 0
\end{align}

\subsubsection{Homogeneous Neumann in 1D, second order}
\begin{align}
	\left\{ \begin{matrix} \ddot p (x,t) = a^2 p''(x,t) \qquad | x \in (0,L) \\ p'(0,t)=0=p'^n_0 \\ p'(L,t)=p'^n_J = \rho d \omega^2 \sin \left( \omega t \right) \qquad | x_J=L \end{matrix} \right.
\end{align}

\begin{align}
	\frac{p^n_1 - p^n_0}{\Delta x} = p'^n_{1/2} + O(\Delta x^2) \overset{\ref{eq:shift1}}{=} p'^n_0 - \frac{\frac{\Delta x}{2}}{1!}p''^n_{1/2}+O(\Delta x^2) \overset{\ref{eq:shift2}}{=} p'^n_0 - \frac{\Delta x}{2}p''^n_{1}+O(\Delta x^2) \overset{\ref{eq:shift3}}{=} \nonumber \\ \overset{\ref{eq:shift3}}{=}  p'^n_0 - \frac{\Delta x}{2}p''^{n-1}_{1}+O(\Delta x^2)\overset{\ref{eq:shift4}}{=} p'^n_0 - \frac{\Delta x}{2c^2\Delta t^2}\left( p^n_1-2p^{n-1}_1+p^{n-2}_1 \right)+O(\Delta x^2)
\end{align}

\begin{align}\label{eq:shift1}
	p'^n_0=p'^n_{1/2} + \frac{\frac{\Delta x}{2}}{1!}p''^n_{1/2}+O(\Delta x^2)
\end{align}

\begin{align}\label{eq:shift2}
	p''^n_{1/2} = p''^n_1 - \frac{\frac{\Delta x}{2}}{1!}p'''^n_1 + O(\Delta x^2) = p''^n_1 + O(\Delta x) 
\end{align}

\begin{align}\label{eq:shift3}
	p''^n_{1} = p''^{n-1}_1 + O(\Delta t) \overset{CFL}{=} p''^{n-1}_1 + O(\Delta x) 
\end{align}

from the equation of motion
\begin{align}\label{eq:shift4}
	p''^{n-1}_{1} = \frac{1}{c^2}\ddot p^{n-1}_1 = \frac{1}{c^2\Delta t^2}\left( p^n_1-2p^{n-1}_1+p^{n-2}_1 \right) + O(\Delta t^2)
\end{align}

\begin{align}
	p^n_0 = p^n_1 - \Delta x p'^n_0 + \frac{\Delta x^2}{2 a^2\Delta t^2 } \left( p^n_1 - 2p^{n-1}_1 + p^{n-2}_1 \right) + O(\Delta x ^3) = \nonumber \\
	= p^n_1 + \frac{\Delta x^2}{2 a^2\Delta t^2 } \left( p^n_1 - 2p^{n-1}_1 + p^{n-2}_1 \right) + O(\Delta x ^3)
\end{align}

\subsubsection{Inhomogeneous Neuann, Transducer in 1D, second order}
similarily as the previous section

\begin{align}
	p^n_J = p^n_{J-1} + \Delta x p'^n_J - \frac{\Delta x^2}{2 a^2\Delta t^2 } \left( p^n_{J-1} - 2p^{n-1}_{J-1} + p^{n-2}_{J-1} \right) + O(\Delta x ^3)
\end{align}

\begin{align}
	p^n_J = p^n_{J-1} + \Delta x \rho d \omega^2 \sin \left( \omega t_n \right) - \frac{\Delta x^2}{2 a^2\Delta t^2 } \left( p^n_{J-1} - 2p^{n-1}_{J-1} + p^{n-2}_{J-1} \right) + O(\Delta x ^3)
\end{align}

$t_n = n\Delta t$

\subsubsection{Impedance or Sommerfeld bc, in 1D, first order}

$\pd{p(x,t)}{t}+Za\pd{p(x,t)}{n}=0$

Note: if the normal of the boundary is $(-1)$ (or in 2D $(-1,0)$) then $n = -x$\\
$a$ is the speed of sound\\
$Z$ is the impedance\\
$Z=0\Rightarrow$ Dirichlet (reflection with same phase)\\
$Z=1\Rightarrow$ Sommerfeld (no reflection, all is absorbed)\\
$Z\rightarrow\infty \Rightarrow$ Neumann (phase inverting reflection)\\
for $n = -x$, with Explicit Euler
\begin{align}
	\frac{p^n_0 - p^{n-1}_0}{\Delta t} = Za\frac{p^{n-1}_1 - p^{n-1}_0}{\Delta x}\\
	p^n_0 = p^{n-1}_0 + \frac{Z a }{\Delta x} \left( p^{n-1}_1 - p^{n-1}_0 \right)
\end{align}
    % ub[0] = ua[0] + Z * dt * c / dx * ( ua[1] - ua[0] ) 

for $n = x$, with Explicit Euler
\begin{align}
	\frac{p^n_J - p^{n-1}_J}{\Delta t} = - Za\frac{p^{n-1}_J - p^{n-1}_{J-1}}{\Delta x}\\
	p^n_J = p^{n-1}_J - \frac{Z a }{\Delta x} \left( p^{n-1}_J - p^{n-1}_{J-1} \right)
\end{align}




\end{document}
